\chapter{Conclusiones}

El desarrollo del Trabajo Práctico N° 5 permitió experimentar y validar métodos indirectos para la caracterización de
inductores y capacitores reales mediante instrumental convencional de laboratorio. A través del análisis de las formas
de onda en el osciloscopio ante excitación triangular, cuadrada y senoidal, se logró cuantificar tanto los parámetros
reactivos concentrados ($L, C$) como sus pérdidas resistivas asociadas ($R$, RSE).

En el método de excitación con corriente triangular sobre la bobina con núcleo de hierro, la adopción de la mediana de los
valores obtenidos ($4.33\text{ H}$) resulta metodológicamente superior a la media aritmética. Esto se debe a que la
presencia del núcleo magnético introduce no linealidades y un marcado crecimiento de las pérdidas magnéticas dinámicas
con la frecuencia, por lo que la mediana actúa como un estimador de tendencia central robusto frente a variaciones
sistemáticas en los extremos de frecuencia. En este mismo ensayo, se confirmó que al invertirse la pendiente de la corriente
triangular ($+m \to -m$), la tensión reactiva autoinducida cambia bruscamente de $+e_L$ a $-e_L$ debido a la condición
de corriente inicial no nula en los flancos de conmutación. En consecuencia, la variación total de tensión observada
de pico a pico es el doble de la amplitud del escalón ($e_{pp} = 2 e_L$), lo que justifica rigurosamente el factor 4 en la
fórmula utilizada.

Asimismo, la señal en bornes del inductor exhibió un aspecto perceptiblemente ruidoso y con sobrepicos transitorios.
Este fenómeno se explica considerando que la inductancia actúa eléctricamente como un circuito derivador
($v_L = L \frac{di}{dt}$), lo cual amplifica cualquier distorsión armónica o ruido de alta frecuencia presente en la
señal de salida del generador.

Por otra parte, la medición de la Resistencia Serie Equivalente (RSE) en capacitores electrolíticos mediante onda
cuadrada demostró ser una técnica directa y eficaz. La cuantificación de la RSE adquiere una relevancia crítica en
aplicaciones prácticas de la ingeniería electrónica. Una RSE elevada degrada el filtrado incrementando el rizado de
tensión a la salida y genera disipación de potencia indeseada por efecto Joule ($P = I_{\text{rms}}^2 \cdot
\text{RSE}$), lo que incrementa la temperatura interna del componente y acorta su vida útil.

En cuanto al método de resonancia LC serie, la determinación experimental de la frecuencia de resonancia ($f_0 = 11.79\si{\kilo\hertz}$)
con un capacitor de referencia permitió calcular la inductancia de la bobina con núcleo de aire ($18.2\mH$) y evaluar su
factor de mérito ($Q = 0.1813$) de forma precisa, identificando además sus frecuencias de corte a $-3\text{ dB}$.
